\documentclass[11pt,a4paper]{article}

\usepackage[utf8]{inputenc}
\usepackage[T1]{fontenc}
\usepackage{amsmath,amssymb,amsthm}
\usepackage[margin=1.1in]{geometry}
\usepackage{xcolor}
\usepackage{hyperref}
\hypersetup{colorlinks=true,linkcolor=blue!50!black,citecolor=black,urlcolor=blue!50!black}

\theoremstyle{definition}
\newtheorem{definition}{Definition}[section]
\theoremstyle{plain}
\newtheorem{theorem}[definition]{Theorem}
\newtheorem{lemma}[definition]{Lemma}
\newtheorem{corollary}[definition]{Corollary}
\theoremstyle{remark}
\newtheorem{remark}[definition]{Remark}

% --- notation ---
\newcommand{\Node}{\mathsf{Node}}
\newcommand{\pa}{\mathrm{pa}}
\newcommand{\dom}{\mathcal{D}}
\newcommand{\Ctx}{\mathcal{U}}         % contexts (exogenous)
\newcommand{\Wld}{\mathcal{W}}         % worlds (endogenous)
\newcommand{\pv}{\Pi}                  % parent view
\newcommand{\solve}{\mathsf{solve}}
\newcommand{\enc}{\mathsf{enc}}
\newcommand{\dect}{\mathsf{dec}}
\newcommand{\elab}{\mathsf{elab}}
\newcommand{\fv}{\mathsf{find\_value}}
\newcommand{\ev}{\mathsf{eval}}
\newcommand{\Dist}[1]{\Delta(#1)}      % distributions on a finite type
\newcommand{\push}[1]{{#1}_{*}}        % pushforward
\newcommand{\Smp}{\mathcal{S}}         % samples
\newcommand{\emp}[1]{\widehat{#1}}
\newcommand{\logv}{\mathsf{log}}
\newcommand{\Ind}[1]{\mathbf{1}_{#1}}
\newcommand{\Ex}{\mathbb{E}}
\newcommand{\Var}{\mathbb{V}}
\renewcommand{\Pr}{\mathrm{Pr}}
\newcommand{\id}{\mathrm{id}}
\newcommand{\dsep}{\perp\!\!\!\perp}
\providecommand{\llbracket}{[\![}
\providecommand{\rrbracket}{]\!]}

\title{\textbf{Experiments as Programs over Causal Models}\\[2mm]
\large A formalised denotational and operational semantics, with
static design-checking as an application}
\author{}
\date{20 October 2026 --- state of the Rocq development}

\begin{document}
\maketitle

\begin{abstract}
We formalise a scientific experiment as a program that transforms a causal
model. The development carries \emph{two} semantics for the same language:
an \emph{abstract} (population-level) semantics, in which a model denotes a
distribution over solved node values and an experiment transforms that
distribution; and a \emph{concrete} (per-unit, finite-sample) semantics, in
which a protocol is run on each of $N$ sampled units and produces a dataset.
The mathematical content of the project is the relation between these two
semantics. Deciding whether an experimental design is \emph{correct} ---
whether its post-experiment graph identifies the target causal query --- is
then a decidable syntactic check, justified by soundness against the
semantics; it is an application of the semantics rather than the object of
study. This document records the definitions and the theorems proved so
far, in mathematical rather than Rocq notation, together with their proofs
and an explicit account of what is \emph{not} yet proved.
\end{abstract}

\tableofcontents

\section{Overview}

The development has three layers.

\begin{enumerate}
\item \textbf{A causal model} is a pair $\langle G, F\rangle$: a DAG $G$
      (the \emph{syntax} --- who may depend on whom) together with a family
      of mechanisms $F$ (the \emph{semantics} --- how), glued by an invariant
      saying that $F$ reads only the parents $G$ declares. Its denotation is
      a distribution over node values, obtained by pushing a law on the
      \emph{unobservables} forward along a deterministic solve. This is
      Pearl's structural-causal-model form: all randomness up front, then
      solve.
\item \textbf{An experiment} is a protocol: a list of operations
      (\emph{intervene}, \emph{randomize}, \emph{measure}) applied to a
      model. It has the two semantics described above.
\item \textbf{A design check} is a decidable predicate on the
      post-experiment graph (e.g.\ the back-door criterion). Each check
      comes with a semantic counterpart and a correspondence theorem.
\end{enumerate}

Everything is finite: all node domains are finite types, so distributions
are finitely-supported and the probability layer is
\texttt{infotheo}'s \texttt{fdist} over an abstract real closed field.
Consequences of that choice, and where it bites, are discussed in
\S\ref{sec:limits}.

Throughout, $\Dist{A}$ denotes the finitely-supported distributions on a
finite type $A$, and $\push{f}\mu$ denotes the pushforward of $\mu$ along
$f$, i.e.\ $(\push{f}\mu)(b) = \sum_{a\,:\,f(a)=b}\mu(a)$.

\section{Preliminaries: graphs, nodes, domains}

\begin{definition}[Graph, nodes]
A graph $G=(V,E)$ has a node list $V$ and an edge list $E$. We write
$\Node(G)$ for the finite type of nodes of $G$ (formally the subtype of
elements of $V$, which makes it a finite type and lets us form dependent
products over it), and $\pa(n)$ for the list of declared parents of $n$,
in edge-list order. We write $N$ for $\Node(G)$ when $G$ is fixed.
\end{definition}

The underlying causal engine (the read-only \texttt{dsep-core} library) is
$\mathbb{N}$-valued: its solver $\fv$ computes with natural numbers. The
typed layer we build on top is \emph{not}; the two are bridged by an
encoding, which is the subject of \S\ref{sec:elab}.

\begin{definition}[Domains, codes]
A \emph{domain family} is a map $\dom : N \to \mathbf{FinType}$, assigning
each node its own finite value space. A \emph{coding} is a pair of families
\[
  \enc_n : \dom_n \to \mathbb{N},
  \qquad
  \dect_n : \mathbb{N} \to \dom_n,
  \qquad
  \dect_n \circ \enc_n = \id .
\]
Note that only this one direction is assumed, and that the reverse cannot
be: $\dect_n$ is total on $\mathbb{N}$ while $\enc_n$'s range is finite, so
$\enc_n(\dect_n(m)) = m$ fails outright for out-of-range $m$. Restricted to
in-range $m$ it holds, but trivially --- if $m = \enc_n(v)$ then
$\dect_n(m) = v$, hence $\enc_n(\dect_n(m)) = m$. Neither version is what
limits Theorem~\ref{thm:solvesource} to parent-free nodes; see
Remark~\ref{rem:depth2} and \S\ref{sec:limits}.
\end{definition}

\begin{remark}[Noise space $=$ value space]
Node $n$'s exogenous slot is taken to have the same type $\dom_n$ as its
value. This is natural because a source node simply takes its unobservable
as its value, so the two spaces must agree there anyway. It is a genuine
restriction: it forbids $|U_n| \neq |\dom_n|$, so principal strata (say a
four-valued compliance type behind a binary outcome) are out of reach until
noise domains are given their own family and codes.
\end{remark}

\section{The causal model}

\subsection{Mechanisms address parents by node}

\begin{definition}[Parent view]
The \emph{parent view} of $n$ is the type
\[
  \pv_n \;=\; \prod_{p \,:\, N} \bigl(\, p \in \pa(n) \;\to\; \dom_p \,\bigr),
\]
a lookup function that answers for exactly those nodes it can be
\emph{proved} are declared parents, returning a value in that parent's own
domain.
\end{definition}

The membership proof is a ticket: no ticket, no read. This is a deliberate
design choice with two consequences. First, parents are addressed by
\emph{node}, never by position, so re-spelling the same DAG (reordering its
edge list) cannot silently change what a mechanism means --- with a
positional interface it can, and the type checker only notices when the
parents' domains happen to differ. Second, ``a mechanism reads only
declared parents'' becomes structural rather than a proof obligation: a
mechanism cannot even \emph{mention} a non-parent, because it could not
supply the ticket.

\begin{definition}[Mechanism]
A \emph{typed mechanism} at $n$ is a function
$F_n : \dom_n \times \pv_n \to \dom_n$, taking $n$'s own noise and a parent
view to $n$'s value. A \emph{typed mechanism family} is $F = (F_n)_{n:N}$.
\end{definition}

\subsection{Contexts and worlds}

The vocabulary is kept strictly separate, because the two spaces mean
different things even though under the simplification above they have the
same shape.

\begin{definition}[Context, world]
A \emph{context} is one experimental unit's complete draw of the
unobservables,
\[
  \Ctx \;=\; \prod_{n\,:\,N} \dom_n ,
\]
and a \emph{world} $\Wld$ is the endogenous outcome: every node's
\emph{solved} value. Both are dependent products over a finite type, hence
finite types, so both can carry distributions.
\end{definition}

All randomness lives in the context and nowhere else; the mechanisms are
deterministic, so a world is a function of a context.

\subsection{Elaboration and the solve}
\label{sec:elab}

\begin{definition}[Elaboration]
$\elab(F)$ is the $\mathbb{N}$-valued mechanism family
\[
  \elab(F)(n)\,(u, \vec{v})
  \;=\;
  \enc_n\Bigl(F_n\bigl(\dect_n(u),\;
      \lambda p.\, \dect_p(\vec v\,[\,\mathrm{idx}(p,\pa(n))\,])\bigr)\Bigr),
\]
where $\mathrm{idx}(p,\pa(n))$ locates $p$ in $n$'s parent list, so a
keyed read is resolved into the positional interface the engine expects.
Nodes outside $G$ keep the engine's default behaviour.
\end{definition}

\begin{lemma}[Elaboration is compatible]
$\elab(F)$ reads only the parents $G$ declares.
\end{lemma}

\begin{proof}
Every read a mechanism can perform is at a declared parent, whose index
into $\pa(n)$ is therefore in range; the engine's positional agreement
hypothesis covers exactly the in-range positions. The ticket discipline of
$\pv_n$ is what makes this hold by construction rather than by inspection
of $F$.
\end{proof}

\begin{definition}[Solve]
Writing $\lceil w \rceil$ for the $\mathbb N$-coded form of a context $w$,
\[
  \solve_F(w)(n) \;=\; \dect_n\bigl(\ev(G,\elab(F),n,\lceil w\rceil)\bigr)
  \;\in\; \dom_n ,
\]
where $\ev$ is the engine's total solver. The whole world determined by $w$
is $\mathsf{world\_of}_F(w) = \lambda n.\,\solve_F(w)(n)$.
\end{definition}

$\ev$ is total by packaging the engine's existence theorem: under
well-formedness and acyclicity, a \emph{complete} exogenous assignment
assigns every node a value. A typed context is complete (every node has a
slot), so the default branch of $\ev$ is provably dead code --- it is not a
clamp, and $\solve$ names an actual solution rather than a fallback.

\section{The abstract (population) semantics}

\subsection{The exogenous law}

\begin{definition}[Exogenous law]
Given a per-node family $P = (P_n)$ with $P_n \in \Dist{\dom_n}$, the joint
law on contexts is the product
\[
  P_\Ctx(w) \;=\; \prod_{n\,:\,N} P_n(w_n).
\]
\end{definition}

The product form \emph{is} the assumption that the exogenous variables are
mutually independent. It is the initial law only: conditioning on an event
about solved values would pull back to a non-product event on contexts, so
after a selection step the unobservables are in general dependent --- that
is not a defect, it is selection bias.

\begin{lemma}[Marginal of the product]
\label{lem:ctxmarg}
$\push{(\lambda w.\,w_n)} P_\Ctx = P_n$.
\end{lemma}

\begin{proof}
Write the indicator $[w_n = v]$ as an extra factor of the product, so that
the mass at $v$ becomes a single sum over contexts of a product over nodes.
Exchange the sum and the product using a \emph{dependent} distributivity
law
\[
  \prod_{i} \sum_{j\,:\,J_i} f(i,j)
  \;=\;
  \sum_{g \,\in\, \prod_i J_i} \prod_i f(i, g(i)),
\]
which is proved by transporting the dependent function space into the
tagged sum $\{i\,\&\,J_i\}$, where the standard non-dependent
distributivity applies. Then every factor except the one at $n$ sums to
$1$, and the factor at $n$ retains the indicator.
\end{proof}

\begin{remark}
\texttt{infotheo} ships only the \emph{homogeneous} product ($N$ copies of
one law, over row vectors). A context is a \emph{heterogeneous} dependent
product, so the dependent distributivity law above had to be proved.
\end{remark}

A query about two nodes at once needs the \emph{joint} marginal, which for a
product law factorises.

\begin{lemma}[Joint marginal at two distinct nodes]
\label{lem:pairmarg}
For $x \neq y$,
$\;\push{(\lambda w.\,(w_x, w_y))} P_\Ctx \;=\; P_x \otimes P_y$.
\end{lemma}

\begin{proof}
As for Lemma~\ref{lem:ctxmarg}, but with \emph{two} indicators inserted as
extra factors, one at $x$ and one at $y$. Exchanging sum and product by the
same dependent distributivity law isolates the factors at $x$ and $y$, each
of which retains its indicator and evaluates to $P_x(v_1)$ resp.\
$P_y(v_2)$, while every other factor sums to $1$. Distinctness of $x$ and
$y$ is what keeps the two isolated factors from colliding.
\end{proof}

\begin{lemma}[A product pushforward is a mixture]
\label{lem:uncurry}
For $F : A \times B \to C$,
\[
  \push{(\lambda (a,b).\,F(a,b))}\,(\mu \otimes \nu)
  \;=\;
  \int_{a \sim \mu} \push{(F(a,\cdot))}\nu ,
\]
i.e.\ pushing a two-argument function through a product law is the same as
drawing the first argument and then pushing the second.
\end{lemma}

\begin{proof}
Evaluate both sides at $c$. The left side is a sum over pairs
$(a,b)$ with $F(a,b) = c$ of $\mu(a)\nu(b)$; re-indexing the sum over
$A \times B$ as an iterated sum over $A$ and then $B$, and pulling the
factor $\mu(a)$ out of the inner sum, gives the right side.
\end{proof}

This is the monadic shape of independence, and it is what turns a
pushforward along a function of two nodes into a mixture over one of them
--- the step that makes randomization visible in
\S\ref{sec:rct2}.

\subsection{The induced laws}

\begin{definition}[World law, node law]
\[
  \llbracket \langle G,F\rangle \rrbracket (P)
  \;=\;
  \push{(\mathsf{world\_of}_F)} P_\Ctx
  \;\in\; \Dist{\Wld},
  \qquad
  \llbracket n \rrbracket (P)
  \;=\;
  \push{(\lambda w.\, \solve_F(w)(n))} P_\Ctx .
\]
We write $\mathsf{fdist\_world}$ and $\mathsf{fdist\_node}(n)$ for these.
\end{definition}

No new randomness is introduced here and none can be: the denotation is a
pushforward of the exogenous law along a deterministic map. The node law
is the abstract program's \emph{per-node output}, and it is what a causal
query ranges over.

\begin{lemma}[Two routes to the node law]
$\push{(\lambda v.\, v_n)}\,\mathsf{fdist\_world} = \mathsf{fdist\_node}(n)$.
\end{lemma}

\begin{proof}
Both sides are pushforwards of $P_\Ctx$; compose the two maps on the left
and observe the composite is pointwise equal to
$\lambda w.\,\solve_F(w)(n)$, then use that pushforward respects pointwise
equality of maps.
\end{proof}

Contrast Lemma~\ref{lem:ctxmarg}: that one recovers a node's \emph{noise}
law and holds by construction, whereas this one has to run the mechanisms.
Only at a source node do the two coincide.

\section{The experiment language}

\begin{definition}[Operations, protocol]
\[
  \mathsf{op} \;::=\;
  \mathsf{Intervene}(n,v) \;\mid\;
  \mathsf{Randomize}(n) \;\mid\;
  \mathsf{Measure}(n),
\]
and a \emph{protocol} is a list of operations. The three span what an
experimenter can do to a causal model and what they can see from it:
intervening and randomizing are \emph{doing} (both cut $n$ from its
parents), measuring is \emph{seeing}. A fourth mode, \emph{selecting}
(inclusion criteria, stratification, loss to follow-up), is deliberately
absent: it is conditioning rather than surgery, it acts on the population
rather than on a unit, and it breaks the product form of the exogenous
law.
\end{definition}

\begin{definition}[Experiment]
An experiment is a record $\langle \text{labels}, F, \text{ops}, N \rangle$:
node labels (marking treatments, responses, and unmeasurable nodes), a
mechanism family, a protocol, and a declared \emph{sample size} $N$.
\end{definition}

The sample size is a hyperparameter of the design, not data: it says how
many units the protocol draws, never which ones. It is recorded because a
trial's sample size is pre-registered --- fixed before any data exists,
exactly like the choice of which node to randomize.

\subsection{Well-formedness, and the ordering discipline}

Well-formedness of an experiment asks that the graph be a well-formed
acyclic DAG with treatments and responses, that every node read be
measurable, that mechanisms be compatible with the DAG (free, by
elaboration), and that no surgery follow a measurement.

That last clause deserves a word. Intervening and randomizing rewrite the
mechanism, so $\mathsf{Measure}(n)$ reads a different value before and
after them; but the abstract layer has no temporal resolution, since
measuring changes neither the graph nor the law. A protocol measuring a
node both before and after surgery would therefore have two concrete
values and only one abstract law. Forbidding surgery after a measurement
rules this out and keeps the two layers in agreement. It is sound but
crude: measuring a node that no later surgery can affect is harmless, so a
descendant-based condition would be strictly more permissive.

\subsection{The concrete per-unit semantics}

\begin{definition}[Per-unit run]
A unit state is a pair (current mechanism family, log), where a log is a
list of (node, value) pairs. One operation acts by
\[
\begin{aligned}
  \mathsf{Intervene}(n,v) &: \quad f_n := \text{const } v, &&\text{log unchanged},\\
  \mathsf{Randomize}(n)   &: \quad f_n := \text{read } U_n, &&\text{log unchanged},\\
  \mathsf{Measure}(n)     &: \quad \text{mechanism unchanged}, &&\text{append } (n, \ev(\dots,n,U)).
\end{aligned}
\]
Running a protocol is the fold of this over the operation list, starting
from the initial mechanism and the empty log.
\end{definition}

Randomization uses a \emph{shadow encoding}: the randomized node reads its
own exogenous slot, and the experimenter's fresh draw is written into that
slot. This matches the mechanism the engine already has, and it fits the
``all randomness up front, then one deterministic solve'' form of the
denotation. The price is that ``randomize equals $\mathrm{do}$'' becomes a
theorem rather than a definition --- Theorem~\ref{thm:randdo}.

\begin{definition}[Logged value]
For an experiment $e$, a context $w$ and a node $n$,
\[
  \logv(e,w,n) \;=\; \dect_n\bigl(\,\text{the entry for } n \text{ in the log of } e \text{ run on } w\,\bigr)
  \;\in\; \dom_n ,
\]
reading an absent entry as the code $0$. Its law across the population is
$\;\mathsf{fdist\_log}(e,n) = \push{(\lambda w.\,\logv(e,w,n))} P_\Ctx$.
\end{definition}

A logged \emph{value} lands in a finite type, where a whole \emph{log} does
not (a log is a list), which is why the per-node reading is what can carry
a distribution. \S\ref{sec:dataset} gives the multi-node version.

\section{The concrete (finite-sample) semantics}

\subsection{Samples must be length-indexed}

\begin{definition}[Sample, sampling law]
A \emph{sample} of size $N$ is a length-indexed vector of contexts,
$\Smp_N = \Ctx^N$, and its law is
\[
  P_{\Smp_N} \;=\; P_\Ctx^{\otimes N}.
\]
\end{definition}

The length must live in the \emph{type}, and this is forced rather than
stylistic. The statement we want --- that an empirical distribution is
close to the true one --- is a claim about the probability that a
\emph{randomly drawn} sample lands close, so the randomness ranges over
samples and samples must carry a law; and a distribution needs a finite
sample space, which a type of lists is not. Consequently there is no single
space of all samples, only a family indexed by $N$: the sample size is
needed in order to \emph{name the space the probability is on}, before the
statement can be written down.

\begin{remark}[Two independence assumptions, now distinct]
$P_{\Smp_N} = \bigl(\bigotimes_n P_n\bigr)^{\otimes N}$ is a product
twice over, and the two products are different claims from different
sources: the inner product over \emph{nodes} is a modelling assumption
about the causal model, the outer product over \emph{units} is a design
assumption about how the sample was drawn (i.i.d.). Naming
$P_{\Smp_N}$ separates them.
\end{remark}

\subsection{The empirical distribution}

\begin{definition}[Empirical distribution]
For $s \in A^{N+1}$,
\[
  \emp{s} \;=\; \push{(\lambda i.\, s_i)}\,\mathrm{Unif}(\{0,\dots,N\})
  \;\in\; \Dist{A},
  \qquad\text{so}\qquad
  \emp{s}(v) = \frac{\#\{\,i : s_i = v\,\}}{N+1}.
\]
\end{definition}

Defining it as a pushforward of the uniform law on \emph{indices}, rather
than by counting, is what makes Lemma~\ref{lem:empmap} immediate.

\begin{lemma}[Tabulating commutes with mapping]
\label{lem:empmap}
For $f : A \to B$ and $s \in A^{N+1}$, \;
$\emp{\langle f(s_i)\rangle_i} = \push{f}\,\emp{s}$.
\end{lemma}

\begin{proof}
Both sides are pushforwards of the uniform law on indices;
$\lambda i.\,f(s_i)$ is the composite of $f$ with $\lambda i.\,s_i$, so
functoriality of the pushforward gives the claim.
\end{proof}

This is the lemma that makes a concrete run comparable with an abstract
pushforward: tabulating a sample and then pushing it through $f$ is the
same as running $f$ on every unit and then tabulating.

\begin{lemma}[The empirical distribution lives on a grid]
\label{lem:grid}
For $s \in A^{N+1}$ and $v \in A$ there is $k \le N+1$ with
$\emp{s}(v) = k/(N+1)$.
\end{lemma}

\begin{proof}
Take $k = \#\{i : s_i = v\}$ and apply the counting form above; $k \le N+1$
because a subset of the index type has at most as many elements as the
index type.
\end{proof}

Lemma~\ref{lem:grid} is the reason an \emph{exact} agreement theorem can
only ever be a warm-up; see the discussion after
Theorem~\ref{thm:exact}.

\subsection{Running the protocol on the whole sample}
\label{sec:dataset}

\begin{definition}[Column, dataset]
For an experiment $e$, a node $n$ and a sample $s \in \Smp_N$, the
\emph{column} at $n$ is
\[
  \mathsf{col}(e,n,s) \;=\; \bigl\langle\, \logv(e, s_i, n) \,\bigr\rangle_{i < N}
  \;\in\; \dom_n^{\,N},
\]
and the \emph{dataset} is the family of columns at the measured nodes,
\[
  \mathsf{data}(e,s) \;=\; \lambda n.\,\lambda (\text{proof that } e \text{ measures } n).\;
  \mathsf{col}(e,n,s).
\]
\end{definition}

The protocol still contains exactly \emph{one} $\mathsf{Measure}(n)$; the
indexed family is what runs that one-operation program at each of the $N$
units. Measuring $N$ units is not expressible by repeating the operation,
since a run consumes one context, so $N$ copies of $\mathsf{Measure}(n)$
would log one unit's value $N$ times. The sample therefore belongs in the
semantics, not in the operation list.

Columns are addressed by node with a ``$e$ measures $n$'' ticket, exactly
as mechanisms address parents: no ticket, no read, so a node the protocol
never measured is \emph{unaskable} rather than answered with the junk value
$\dect_n(0)$ that an absent log entry would produce.

What having all columns together buys is \emph{joint} observation: unit
$i$'s pair of readings at two measured nodes $x,y$ is
$\bigl(\mathsf{data}(e,s)(x)_i,\ \mathsf{data}(e,s)(y)_i\bigr)$, indexed by
a genuine $i < N$, so the pairing cannot drift. A per-node marginal cannot
express a joint, and a joint is what a query such as $P(R \mid do(T))$
requires.

\begin{remark}
A ticketed family is a function out of a proposition, hence not a finite
type, so the dataset cannot itself carry a distribution. This costs
nothing here: a dataset is a \emph{deterministic function of the sample},
and every probabilistic statement below is phrased over $P_{\Smp_N}$. The
concrete side holds data, the abstract side holds the law, and the two are
compared through $\emp{\cdot}$ --- so the law lives where the randomness
is, and the dataset is free to be faithful rather than enumerable.
\end{remark}

\section{Theorems}

\subsection{Elaboration correctness at a source node}

\begin{theorem}[Parent-free nodes solve to what their mechanism says]
\label{thm:solvesource}
Let $G$ be well-formed and acyclic, let $\pa(n) = \varnothing$, and suppose
$n$'s mechanism is \emph{parent-free}: $F_n(u,\pi) = g(u)$ for some
$g : \dom_n \to \dom_n$ and all $u,\pi$. Then for every context $w$,
\[
  \solve_F(w)(n) \;=\; g(w_n).
\]
\end{theorem}

\begin{proof}
By the engine's existence theorem, under well-formedness, acyclicity and a
complete exogenous assignment, $\fv$ computes $n$'s value from its own
slot and its parents' values. At $n$ the parent list is empty, and reading
$n$'s slot out of the coded context returns exactly $\enc_n(w_n)$ (the
read-after-write property of the coding bridge, which holds because node
identifiers are distinct). Unfolding $\elab$ at $n$ gives
\[
  \enc_n\bigl(F_n(\dect_n(\enc_n(w_n)),\,\cdot\,)\bigr)
  = \enc_n\bigl(F_n(w_n, \cdot)\bigr)
  = \enc_n(g(w_n)),
\]
using $\dect_n \circ \enc_n = \id$ and then the hypothesis on $F_n$.
Finally $\solve$ applies $\dect_n$ to this, and $\dect_n\circ\enc_n = \id$
again.
\end{proof}

\begin{remark}
This is the first proved instance of elaboration correctness --- the first
place where the solver is shown to compute what the typed mechanism
\emph{says}, rather than merely something total. It works because the value
round-trips $\enc \to \dect \to \enc \to \dect$, so every step is the
available cancellation; it never needs range closure
$\enc_n(\dect_n(m)) = m$, which is what blocks the general statement.

Note also that $g$ is not a restriction. With $\pa(n)=\varnothing$ the type
$\pv_n$ has essentially one inhabitant, so \emph{every} mechanism at $n$ is
already parent-independent; $g$ merely has to be named in order to say what
the answer is.
\end{remark}

\subsection{The experiment-level law}

\begin{theorem}[Law of the logged value]
\label{thm:law}
Under the hypotheses of Theorem~\ref{thm:solvesource}, for an experiment
$e$ whose protocol is exactly $[\mathsf{Measure}(n)]$,
\[
  \mathsf{fdist\_log}(e,n) \;=\; \push{g}\,P_n .
\]
\end{theorem}

\begin{proof}
Because the protocol consists of the single operation
$\mathsf{Measure}(n)$, the log has one entry, keyed by $n$, and reading it
back gives $\solve_F(w)(n)$; hence
$\logv(e,w,n) = \solve_F(w)(n) = g(w_n)$ by
Theorem~\ref{thm:solvesource}. So
\[
  \mathsf{fdist\_log}(e,n)
  = \push{(\lambda w.\,g(w_n))} P_\Ctx
  = \push{g}\,\push{(\lambda w.\,w_n)} P_\Ctx
  = \push{g}\, P_n,
\]
by functoriality of the pushforward and then Lemma~\ref{lem:ctxmarg}.
\end{proof}

\begin{corollary}[``Hello world'': a pure echo logs the noise law]
If $F_n(u,\pi) = u$ then $\mathsf{fdist\_log}(e,n) = P_n$.
\end{corollary}

\begin{proof}
The case $g = \id$ of Theorem~\ref{thm:law}, with
$\push{\id}\mu = \mu$.
\end{proof}

Theorem~\ref{thm:law} is the end-to-end smoke test of the entire stack ---
coding, the context bridge, the engine's solver, the total evaluator, the
per-unit runner, the product law --- and it is an \emph{exact equality of
laws}, not a limit statement. It is also the first fragment of the general
pushforward story: a parent-free node already exhibits the shape ``law of
the logged value $=$ pushforward of the exogenous law along what the
mechanism does''.

\subsection{Exact agreement on a representative sample}

\begin{theorem}[Exact-sample law]
\label{thm:exact}
Under the hypotheses of Theorem~\ref{thm:law}, for a sample
$s \in \Smp_{N+1}$ that is \emph{representative}, i.e.\
$\emp{s} = P_\Ctx$,
\[
  \emp{\mathsf{col}(e,n,s)} \;=\; \push{g}\, P_n .
\]
\end{theorem}

\begin{proof}
The column is $\langle \logv(e,s_i,n)\rangle_i$, so by
Lemma~\ref{lem:empmap}
\[
  \emp{\mathsf{col}(e,n,s)}
  = \push{(\lambda w.\,\logv(e,w,n))}\,\emp{s}
  = \push{(\lambda w.\,\logv(e,w,n))}\,P_\Ctx,
\]
using representativeness, and the right-hand side is
$\mathsf{fdist\_log}(e,n) = \push{g}P_n$ by Theorem~\ref{thm:law}.
\end{proof}

This is the honest form of ``exhaustive'': the sample is assumed to carry
the right law, and the assumption is \emph{stated} rather than assumed
away. Everything in sight is finite --- there are no limits anywhere.

\begin{remark}[Why this is a warm-up and not the result]
\label{rem:warmup}
The representativeness premise silently restricts $P$. By
Lemma~\ref{lem:grid} an empirical distribution takes only values
$k/(N+1)$, so $\emp{s} = P_\Ctx$ is satisfiable only when every value of
$P_\Ctx$ is rational with denominator dividing $N+1$. Against an
irrational truth the theorem is vacuous. Moreover no choice of number
field repairs this: even with rational probabilities a real sample is not
exactly representative --- drawing $400$ units from a population with
$p = 1/3$ cannot give exactly $400/3$ successes. For irrational $p$ the
error is bounded \emph{below} by $p$'s distance to the grid, which is
positive. An $\varepsilon$-bound is therefore not a weakening adopted for
convenience; it is the only shape of theorem available. That is
\S\ref{sec:lln}.
\end{remark}

\subsection{A weak law of large numbers for the outcome}
\label{sec:lln}

Fix a value $v \in \dom_n$ and let $A_v \subseteq \Ctx$ be the event that a
unit logs $v$,
$A_v = \{\,w : \logv(e,w,n) = v\,\}$, with indicator $\Ind{A_v}$.

\begin{lemma}[The indicator's mean is the outcome law]
\label{lem:eind}
Under the hypotheses of Theorem~\ref{thm:law},
$\;\Ex_{P_\Ctx}\bigl[\Ind{A_v}\bigr] = (\push{g}P_n)(v)$.
\end{lemma}

\begin{proof}
The mean of an indicator is the probability of its event, and
$P_\Ctx(A_v)$ is by definition the mass that
$\push{(\lambda w.\,\logv(e,w,n))}P_\Ctx = \mathsf{fdist\_log}(e,n)$
assigns to $v$; apply Theorem~\ref{thm:law}.
\end{proof}

\begin{lemma}[The empirical frequency is the averaged indicator]
\label{lem:bridge}
For $f : A \to B$, $v \in B$ and $s \in A^{N+1}$,
\[
  \emp{\langle f(s_i)\rangle_i}(v)
  \;=\;
  \frac{1}{N+1}\sum_{i<N+1} \Ind{\{a\,:\,f(a)=v\}}(s_i).
\]
\end{lemma}

\begin{proof}
Both sides count the indices at which $f(s_i) = v$ and divide by $N+1$: the
left by the counting form of $\emp{\cdot}$, the right by splitting the sum
on the predicate $f(s_i) = v$, where the complementary part contributes
zero.
\end{proof}

Lemma~\ref{lem:bridge} is the bridge that makes a statement about the
program's output a statement the weak law can act on, with no new
probability theory required.

\begin{theorem}[Weak law for the measured outcome]
\label{thm:lln}
Under the hypotheses of Theorem~\ref{thm:law}, if the experiment declares
sample size $N+1$ and $\varepsilon > 0$, then
\[
  \Pr_{\,s \,\sim\, P_{\Smp_{N+1}}}
  \Bigl[\;
     \bigl|\;\emp{\mathsf{col}(e,n,s)}(v) \;-\; (\push{g}P_n)(v)\;\bigr|
     \;\ge\; \varepsilon
  \;\Bigr]
  \;\le\;
  \frac{\Var_{P_\Ctx}\bigl[\Ind{A_v}\bigr]}{(N+1)\,\varepsilon^{2}} .
\]
\end{theorem}

\begin{proof}
By Lemma~\ref{lem:bridge} the quantity inside the probability is the
average of $N+1$ independent copies of $\Ind{A_v}$, minus its mean --- the
mean being $(\push{g}P_n)(v)$ by Lemma~\ref{lem:eind}. Chebyshev's
inequality applied to that average, whose variance is
$\Var[\Ind{A_v}]/(N+1)$, gives the bound. Formally this instantiates the
library's weak law, whose hypotheses are that the summands are i.i.d.\
copies with common mean and common variance; the sum-decomposition
obligation is discharged by a generic ``sum of $n$ i.i.d.\ copies of an
arbitrary random variable'' construction, which the library did not
provide (it instantiates the relation only for one specific variable in its
asymptotic-equipartition development).
\end{proof}

Three things are worth emphasising about the shape of
Theorem~\ref{thm:lln}.

\begin{itemize}
\item \textbf{It is a finite-$N$ bound, not a limit.} This is what makes it
      provable in this setting at all: distributions here are
      finitely-supported, so the infinite product space that almost-sure
      convergence requires is not expressible. A $1/N$ bound delivers
      convergence in probability without ever mentioning a limit.
\item \textbf{It carries no representativeness premise}, so unlike
      Theorem~\ref{thm:exact} it is non-vacuous for arbitrary real-valued
      $P$ (cf.\ Remark~\ref{rem:warmup}). The grid of
      Lemma~\ref{lem:grid} has spacing $1/(N+1) \to 0$ while $\varepsilon$
      stays positive, so there is always room; convergence never requires
      hitting the limit.
\item \textbf{The declared sample size sets the rate.} The bound quotes the
      size recorded in the experiment, so the hypothesis tying it to $N+1$
      is load-bearing: without it the statement would quote a size with no
      relation to the sample at hand.
\end{itemize}

\subsection{The checker, and its soundness}

The static design checks are the application layer. Each consists of a
decidable predicate on the post-experiment graph, a semantic counterpart,
and a correspondence theorem; what follows is the state of that layer.

\begin{theorem}[Randomization is a $\mathrm{do}$, per unit]
\label{thm:randdo}
Let $G$ be well-formed and acyclic, let $n,w$ be nodes of $G$, and let $U$
be a complete exogenous assignment. Then
\[
  \fv\bigl(G,\ \text{randomize } n \text{ in } F,\ w,\ (n\mapsto v)\cdot U\bigr)
  \;=\;
  \fv\bigl(\mathrm{do}(n)\,G,\ F[f_n := \text{const } v],\ w,\ U\bigr).
\]
\end{theorem}

That is, the shadow encoding (write the draw into $n$'s exogenous slot and
let $f_n$ read it) and the surgery encoding (cut $n$'s incoming edges and
fix $f_n$ to the constant) compute the same value at every node. This is
proved by induction along a topological order. It is the per-unit core of
the adequacy theorem: the distributional statement is obtained by lifting
it through the mixture over the coin, which is the next milestone.

\begin{theorem}[Back-door correspondence]
Under the standing graph hypotheses, with $R$ a node of $G$, $T \neq R$,
$Z \subseteq V$ and $T, R \notin Z$,
\[
  \mathsf{syntactic\_backdoor}(G,T,R,Z) = \mathrm{true}
  \iff
  \mathsf{semantic\_backdoor}(G,T,R,Z),
\]
where the syntactic side asks that no element of $Z$ be a descendant of $T$
and that $T$ and $R$ be $d$-separated by $Z$ in $G$ with $T$'s outgoing
edges removed, and the semantic side replaces $d$-separation by its
semantic counterpart.
\end{theorem}

This is the per-mechanism instance of the soundness direction: a decidable
graph condition implies the corresponding independence. It is obtained from
the underlying library's equivalence of $d$-separation and semantic
separation; the conditioning set is deduplicated internally so that
duplicate-freeness need not be assumed.

\begin{theorem}[A randomized trial is syntactically correct]
For the protocol ``randomize $T$, then measure $R$'', the post-experiment
DAG is $\mathrm{do}(T)\,G$, in which $T$ is \emph{isolated}; hence the
back-door criterion holds with $Z = \varnothing$.
\end{theorem}

\begin{proof}
Randomizing $T$ removes $T$'s incoming edges, and the back-door
mutilation removes its outgoing edges; so in the mutilated graph $T$ has no
incident edges at all, and $d$-separation from $R$ with an empty
conditioning set is immediate.
\end{proof}

The isolation argument is what makes this short. An earlier formulation
introduced a fresh randomizer node feeding $T$, and the corresponding proof
required an extensive adjacency analysis; keeping the node set fixed and
reading randomization as surgery on the existing node replaced that with
three small lemmas.

\section{A worked experiment: the two-node randomized trial}
\label{sec:rct2}

Theorem~\ref{thm:law} measures one parentless node. This section does the
smallest thing that is genuinely an experiment, and two features are new:
the measured node \emph{has a parent}, and the conclusion is a
\emph{mixture}.

\subsection{The model and the protocol}

Let $G_2$ have nodes $T, R$ and the single edge $T \to R$, with
\[
  \dom_T = \{\textsf{placebo}, \textsf{drug}\},
  \qquad
  \dom_R = \{\textsf{worse}, \textsf{same}, \textsf{better}, \textsf{cured}\},
\]
and mechanisms, for an arbitrary \emph{response table}
$h : \dom_T \to \dom_R \to \dom_R$,
\[
  F_T(u,\pi) \;=\; \textsf{placebo},
  \qquad
  F_R(u,\pi) \;=\; h\bigl(\pi(T),\, u\bigr).
\]
$F_T$ is deliberately \emph{biased}: it ignores its noise and assigns
everyone placebo. It stands for a natural assignment mechanism --- a
clinician's judgement, say --- that the trial exists to override. The
protocol is
\[
  \bigl[\;\mathsf{Randomize}(T),\ \mathsf{Measure}(R)\;\bigr].
\]
Write $F'$ for the post-randomization family, i.e.\ $F$ with $f_T$ replaced
by ``return your own exogenous slot''.

\subsection{The solve, at depth two}

\begin{lemma}[$T$ after randomization]
\label{lem:rct2T}
$\fv(G_2, F', T, \lceil w \rceil) = \enc_T(w_T)$.
\end{lemma}

\begin{proof}
$T$ is parentless, so the engine evaluates $f'_T$ on $T$'s own slot and an
empty parent list. After randomization $f'_T$ \emph{is} ``return the
noise'', so it hands back the stored code unchanged. No decoding happens at
all here.
\end{proof}

\begin{lemma}[$R$: the induction step]
\label{lem:rct2R}
$\fv(G_2, F', R, \lceil w \rceil) = \enc_R\bigl(h(w_T, w_R)\bigr)$.
\end{lemma}

\begin{proof}
By the engine's characterisation, $R$'s value is $f'_R$ applied to $R$'s own
slot together with the list of its parents' computed values. $R$'s parent
list is $[T]$, whose single entry is $\enc_T(w_T)$ by
Lemma~\ref{lem:rct2T}. Unfolding the elaboration at $R$ gives
\[
  \enc_R\Bigl(F_R\bigl(\dect_R(\enc_R(w_R)),\;
     \lambda p.\,\dect_p(\,\cdot\,)\bigr)\Bigr),
\]
where the parent lookup resolves $T$'s index in $[T]$ to position $0$ and
so decodes $\enc_T(w_T)$ to $w_T$. Both decodings --- $R$'s own noise and
$T$'s parent value --- are instances of $\dect \circ \enc = \id$, so the
value is $\enc_R(h(w_T, w_R))$.
\end{proof}

\begin{remark}[Why this matters beyond the example]
\label{rem:depth2}
Lemma~\ref{lem:rct2R} is the depth-two case of the general
elaboration-correctness statement
\begin{equation}
  \solve_F(w)(n) \;=\; F_n\bigl(w_n,\ \lambda p.\,\solve_F(w)(p)\bigr),
  \label{eq:elabcorrect}
\end{equation}
which \S\ref{sec:limits} lists as open, and it shows where the difficulty
is \emph{not}.

It is not in the coding. One might expect \eqref{eq:elabcorrect} to need
the reverse cancellation $\enc(\dect(m)) = m$, but it needs no such thing.
The parent view that $\elab$ supplies hands the mechanism
$\dect_p(\text{$p$'s slot in the positional list})$, while $\solve_F(w)(p)$
is \emph{by definition} $\dect_p$ of the solver's value at $p$. So once that
slot is known to hold the solver's value at $p$, the two sides agree
definitionally, and no claim about which naturals are in range appears
anywhere. All the arithmetic is the one available cancellation, used twice:
on $n$'s own noise, and to strip the outer $\enc_n$. That is precisely why
the proof of Lemma~\ref{lem:rct2R} consists of applying it twice and
nothing else.

What \eqref{eq:elabcorrect} does need is a \textbf{positional--keyed
correspondence}. The engine hands a mechanism its parents' values as a
\emph{list, in the order of $\pa(n)$}; $\elab$ reads a parent \emph{by
node}, resolving it through that parent's index into the same list. Nothing
so far proves the two agree, and that --- bookkeeping about the engine's
internals, not about the coding --- is the open obligation. In the present
example the correspondence holds by computation, because $\pa(R) = [T]$ and
the index is $0$.

Note finally that no induction on the graph seems to be required: the
engine's characterisation already unfolds one step at \emph{every} node
independently, and \eqref{eq:elabcorrect} only needs $\solve_F(w)(p)$ to be
\emph{named} on the right-hand side, not computed.
\end{remark}

\subsection{The outcome law}

\begin{theorem}[Law of the logged response]
\label{thm:rct2law}
$\;\mathsf{fdist\_log}(e, R) = \push{(\lambda w.\, h(w_T, w_R))} P_\Ctx$.
\end{theorem}

\begin{proof}
The protocol logs one entry, at $R$; reading it back gives
$\dect_R$ of Lemma~\ref{lem:rct2R}, i.e.\ $h(w_T, w_R)$. The law is then
the pushforward of $P_\Ctx$ along that map, by definition.
\end{proof}

This is true but not yet informative: it still mentions the joint context,
and the same statement holds for \emph{any} protocol that measures $R$.
What randomization buys is that it factorises.

\begin{theorem}[Correctness of the trial]
\label{thm:rct2}
\[
  \mathsf{fdist\_log}(e, R)
  \;=\;
  \int_{t \sim P_T} \push{\bigl(h(t,\cdot)\bigr)} P_R .
\]
That is: the outcome law is the coin-\textbf{mixture} of the per-arm
response laws --- draw the treatment from $T$'s law, then push $R$'s noise
through the response table at that treatment.
\end{theorem}

\begin{proof}
By Theorem~\ref{thm:rct2law} it suffices to factorise the pushforward. The
map $w \mapsto h(w_T, w_R)$ is the composite of $w \mapsto (w_T, w_R)$ with
the uncurried $h$. The first pushes $P_\Ctx$ forward to $P_T \otimes P_R$
by Lemma~\ref{lem:pairmarg}, since $T \neq R$; and pushing the uncurried
$h$ through a product law is exactly the mixture, by
Lemma~\ref{lem:uncurry}.
\end{proof}

\subsection{What the example shows}

\begin{itemize}
\item \textbf{Randomization erased the biased mechanism.} $F_T$ assigns
      everyone placebo, and it does not appear in
      Theorem~\ref{thm:rct2}'s conclusion. The surgery really did
      something, and the theorem is the witness.
\item \textbf{The mixture is the content.} Independence of $T$'s value from
      $R$'s noise is what makes the outcome law decompose by treatment arm,
      and that decomposition is what an experimenter relies on when
      comparing arms. Theorem~\ref{thm:rct2law} alone would not support
      that reading.
\item \textbf{The coin needs no new machinery here.}
      Theorem~\ref{thm:rct2} is general in $P$, so ``the coin at $T$'' is
      simply which $P$ one instantiates it at; a balanced trial is
      $P_T$ uniform, an unequal allocation is any other $P_T$, and nothing
      in the proof depends on the choice. A \emph{law-level} randomization
      operation --- replacing $P_T$ by a declared coin --- becomes
      necessary only for a compositional runner that walks the operation
      list. Note also that hoisting such law-surgery to before sampling is
      sound only because the ordering discipline forces all surgery to
      precede all measurement; it would not survive adaptive designs.
\item \textbf{It is still one node's marginal.} The theorem is about $R$'s
      law, not about the joint law of $(T,R)$ that a conditional query would
      need. Lemma~\ref{lem:pairmarg} is the piece that makes such joints
      expressible; using it to state an identification result is Theorem B,
      which remains open.
\end{itemize}

\section{Confounding, and what randomization removes}
\label{sec:rct3}

The trial of \S\ref{sec:rct2} overrides a treatment mechanism that was
merely \emph{biased}. This section does the case randomization actually
exists for --- a \textbf{confounder} --- and proves the two protocols
denote different laws, so that the intervention is not a no-op.

\subsection{The model}

Let $G_3$ have nodes $X$ (treatment), $Y$ (outcome) and $Z$ (a common
cause), with edges $X \to Y$, $Z \to X$, $Z \to Y$:
\[
\begin{array}{ccc}
     & Z & \\
 \swarrow & & \searrow \\
 X & \longrightarrow & Y
\end{array}
\]
Read as a small drug trial: $Z$ is age with
$\dom_Z = \{\textsf{young},\textsf{middle},\textsf{old}\}$,
$\dom_X = \{\textsf{placebo},\textsf{drug}\}$, and
$\dom_Y = \{\textsf{worse},\textsf{same},\textsf{better},\textsf{cured}\}$
--- three categorical domains of three different cardinalities. The
mechanisms are
\[
  F_Z(u,\pi) = u,
  \qquad
  F_X(u,\pi) = \beta\bigl(\pi(Z)\bigr),
  \qquad
  F_Y(u,\pi) = h\bigl(\pi(X),\,\pi(Z)\bigr),
\]
where $\beta(\textsf{old}) = \textsf{drug}$ and $\beta = \textsf{placebo}$
otherwise --- ``the old get treated'' --- and $h$ is the response table
\[
\begin{array}{llll}
  h(\textsf{drug},\textsf{young}) = \textsf{cured}, &
  h(\textsf{drug},\cdot) = \textsf{better}, &
  h(\textsf{placebo},\textsf{old}) = \textsf{worse}, &
  h(\textsf{placebo},\cdot) = \textsf{same}.
\end{array}
\]
$F_X$ is the confounding: the treatment \emph{reads} the common cause, so
in the untouched model $X$ and $Z$ are dependent. Note that $Y$ now has
\emph{two} parents, whose positional list is $[X; Z]$, so the keyed reads
$\pi(X)$ and $\pi(Z)$ must resolve to positions $0$ and $1$ --- the first
place where the lookup discipline of \S\ref{sec:elab} is genuinely
exercised, since in \S\ref{sec:rct2} every parent list had one element.

We compare two protocols on this one model:
\[
  e_{\mathrm{trial}} = \bigl[\,\mathsf{Randomize}(X),\ \mathsf{Measure}(Y)\,\bigr],
  \qquad
  e_{\mathrm{obs}} = \bigl[\,\mathsf{Measure}(Y)\,\bigr].
\]

\subsection{The trial}

\begin{lemma}
\label{lem:rct3}
Under $e_{\mathrm{trial}}$, $\;\solve(w)(Y) = h(w_X, w_Z)$.
\end{lemma}

\begin{proof}
$Z$ is a source and echoes its noise, giving $w_Z$. At $X$ the mechanism
has been replaced by ``return your own slot'', so $X$ takes $w_X$. The
engine still \emph{builds and passes} $X$'s parent list $[Z]$, and the
replaced mechanism simply ignores it --- that discarded list is the edge
$Z \to X$ being severed. At $Y$ the positional list is
$[\enc_X(w_X), \enc_Z(w_Z)]$, and the keyed reads recover $w_X$ and $w_Z$
by the cancellation, so $Y$'s value is $h(w_X,w_Z)$.
\end{proof}

\begin{theorem}[Outcome law of the trial]
\label{thm:rct3}
\[
  \mathsf{fdist\_log}(e_{\mathrm{trial}}, Y)
  \;=\;
  \int_{t \sim P_X} \push{\bigl(h(t,\cdot)\bigr)} P_Z .
\]
\end{theorem}

\begin{proof}
As for Theorem~\ref{thm:rct2}: push Lemma~\ref{lem:rct3} forward, factorise
the joint at the distinct nodes $X$ and $Z$ by Lemma~\ref{lem:pairmarg},
and convert the product pushforward into a mixture by
Lemma~\ref{lem:uncurry}.
\end{proof}

The confounder has not left the model --- $h$ still reads it, so age still
moves the outcome. What has gone is the \emph{correlation} between age and
treatment, and the factorisation is the formal trace of its removal: the
outer integral ranges over a treatment law that no longer depends on $Z$.

\subsection{The same measurement without randomizing}

\begin{lemma}
\label{lem:obs3}
Under $e_{\mathrm{obs}}$, $\;\solve(w)(Y) = h\bigl(\beta(w_Z),\, w_Z\bigr)$.
\end{lemma}

\begin{proof}
Now $F_X$ is live, so the chain is three deep: $Z$ echoes $w_Z$; $X$ reads
$Z$ through its parent list and returns $\beta(w_Z)$; $Y$ reads both and
returns $h(\beta(w_Z), w_Z)$. Each parent value arrived as $\enc$ of
something and is recovered by the cancellation, exactly as at depth two.
\end{proof}

\begin{theorem}[Observational law]
\label{thm:obs3}
\[
  \mathsf{fdist\_log}(e_{\mathrm{obs}}, Y)
  \;=\;
  \push{\bigl(\lambda a.\, h(\beta(a), a)\bigr)} P_Z .
\]
\end{theorem}

\begin{proof}
Push Lemma~\ref{lem:obs3} forward. The resulting map factors through
$w \mapsto w_Z$, so the single-node marginal Lemma~\ref{lem:ctxmarg}
replaces $P_\Ctx$ by $P_Z$.
\end{proof}

This law is a function of the \textbf{confounder alone}: $Z$ fixes the
treatment, which fixes the response. There is no treatment law left to
integrate over, and that absence is precisely the sense in which the design
is biased --- not an error in the calculation, but a missing degree of
freedom in the design.

\subsection{The two differ}

\begin{theorem}[Randomizing is not a no-op]
\label{thm:sep3}
At the concrete exogenous law $P_3$ --- $Z$ uniform on the three ages ---
\[
  \mathsf{fdist\_log}(e_{\mathrm{trial}}, Y)
  \;\neq\;
  \mathsf{fdist\_log}(e_{\mathrm{obs}}, Y),
\]
witnessed at the outcome $\textsf{worse}$, which the trial gives mass
$\tfrac{1}{3}$ and the observational protocol mass $0$.
\end{theorem}

\begin{proof}
The only pair with $h(t,a) = \textsf{worse}$ is
$(\textsf{placebo}, \textsf{old})$. Under Theorem~\ref{thm:rct3} the
treatment and the age are independent, so that pair has probability
$P_X(\textsf{placebo}) \cdot P_Z(\textsf{old}) = \tfrac{1}{3}$ at $P_3$.
Under Theorem~\ref{thm:obs3} the outcome is $\textsf{worse}$ only if
$\beta(a) = \textsf{placebo}$ and $a = \textsf{old}$; but
$\beta(\textsf{old}) = \textsf{drug}$, so no age qualifies and the mass is
$0$. Since $\tfrac13 \neq 0$ the two laws are distinct.
\end{proof}

The mechanism of the separation is worth stating in words: $\textsf{worse}$
arises only from \emph{placebo and old}, and the confounded mechanism
treats exactly the old --- so under $e_{\mathrm{obs}}$ no unit can ever be
both, and the outcome is unreachable. Randomization decouples the two and
the mass appears.

\begin{remark}[An honest caveat about this instance]
At $P_3$ the exogenous law at $X$'s slot is a point mass on
$\textsf{placebo}$, inherited from the original model where $X$'s mechanism
ignored its noise. So the particular trial separated here is
\emph{degenerate}: after randomization every unit receives placebo.
Theorem~\ref{thm:sep3} is still a genuine separation --- and the
degeneracy sharpens it, since under $e_{\mathrm{obs}}$ the old receive the
drug and reach $\textsf{better}$ instead --- but it is not an illustration
of a balanced trial. Theorem~\ref{thm:rct3} is general in $P$, so a
balanced allocation is simply a different instantiation; nothing in its
proof depends on the coin.
\end{remark}

\begin{remark}[Further evidence for Remark~\ref{rem:depth2}]
Lemma~\ref{lem:obs3} is a \emph{three}-deep solve, with $X$ genuinely
reading $Z$, and like Lemma~\ref{lem:rct2R} its entire arithmetic is the
one available cancellation applied pointwise. No induction on the graph
appears, and no claim about which naturals lie in $\enc$'s range. Together
with the two-parent lookup at $Y$, this is the strongest available evidence
that the general statement \eqref{eq:elabcorrect} needs only the
positional--keyed correspondence.
\end{remark}

\section{What is not proved}
\label{sec:limits}

Stated plainly, so that the scope of the above is not overread.

\begin{description}
\item[General elaboration correctness.] Theorem~\ref{thm:solvesource}
  covers parent-free nodes, and Lemma~\ref{lem:rct2R} one node with one
  parent; the general statement, that the solver computes what a typed
  mechanism means at \emph{every} node, is not proved. Until it is,
  $\solve$ at a node with parents is ``whatever the engine computes,
  decoded'' rather than provably ``what the mechanism says''. The
  outstanding obligation is the positional--keyed correspondence between
  the engine's parent-value list and $\elab$'s keyed reads, \emph{not} any
  property of the coding --- see Remark~\ref{rem:depth2}.

\item[Adequacy (Theorem A).] Theorem~\ref{thm:randdo} is the per-unit
  core; the distributional lift --- that the compositional population
  semantics agrees with ``run each unit and push forward'' --- is not yet
  proved. This is the main outstanding result, and it is what transfers
  the believability of the small operational semantics to the
  reasoning-friendly denotational one.

\item[Identification (Theorem B, distributional).] That the log of a
  randomized trial identifies $P(R \mid do(T))$ is not yet proved. It needs
  the coin to have full support (overlap) and to be independent of
  confounders --- both hold for a randomized trial by construction, and
  uniformity is \emph{not} required.

\item[The weak law is per-value.] Theorem~\ref{thm:lln} fixes one
  $v \in \dom_n$. A statement about the whole distribution needs a union
  bound, at a cost of a factor $|\dom_n|$.

\item[The variance is left abstract.] $\Var[\Ind{A_v}] = p(1-p) \le 1/4$ is
  not computed, so the bound is stated in terms of a variance rather than
  a constant.

\item[Nothing downstream of the context bridge computes.] The bridge
  enumerates the node type, and that enumeration does not reduce, so
  per-node laws can be \emph{stated} but not \emph{evaluated}. This is not
  a defect for a specification layer, but it does mean the development is
  not a simulator.

\item[Selection is absent.] Conditioning on measured values would require
  positive mass --- a side condition that well-formedness cannot discharge,
  since it depends on the numbers --- and it destroys the product form of
  the exogenous law. The design was worked out and deliberately deferred.

\item[Continuous outcomes.] Finiteness of the sample space is what the
  whole probability layer rests on. A real-valued node must be discretised,
  or the development must move to a measure-theoretic setting.

\item[Scope of the language.] Treatments and responses are single nodes,
  not vectors; the protocol is a straight-line list, so adaptive designs
  and dynamic treatment regimes (which need branching on measured data)
  are not expressible.
\end{description}

\section{Metatheory}

Every theorem above is machine-checked in Rocq 9.1.0 against
\texttt{mathcomp} 2.5.0 and \texttt{infotheo} 0.9.7, with no admitted
statements. The axiom footprint is exactly the three classical axioms the
underlying analysis library already assumes --- propositional
extensionality, dependent functional extensionality, and constructive
indefinite description --- and the development adds none of its own.

\end{document}
